\section{Introduction}\label{sec:intro}

A phase-only SLM placed in the pupil plane of a lens forms an array of optical traps whose number now reaches the thousands. The largest
published array holds about 6100 atoms in about 12,000 tweezers, and the same report states that the light field is one of the limits of
the machine: the trap depth varies by 11.4 percent (standard deviation) across the sites after weighted GS optimization, the finite number of SLM pixels
bounds the array, and aberrations grow with the tweezer power~\cite{manetsch}. A small calibrated array reaches a relative standard deviation
of the intensity of 0.3 percent~\cite{chew}, so the spread grows with the size of the array and with the distance between the model that computes the hologram
and the device that displays it. Error-corrected neutral-atom processors of hundreds of physical qubits~\cite{bluvstein} and hybrid architectures that keep more than 3000 atoms for hours with a
static SLM array~\cite{chiu} make the quality of the hologram a question of the experiment's yield, not of its elegance.

The hologram is computed, not designed, and one family of method dominates large arrays. The Gerchberg--Saxton iteration and its weighted form~\cite{dileonardo}
equalize the trap intensities in the ideal model of the SLM; freezing the focal-plane phase after a few iterations (phase-fixed weighted GS)~\cite{kim2019} keeps the weights effective and reaches
percent-level uniformity on arrays of 1500 foci; the same algorithm computes the holograms of the arrays of about 1000 tweezers and of 11,998 sites~\cite{ebadi,manetsch}. Camera feedback then
corrects the residual spread on the experiment~\cite{nogrette,kim2019}. The solver itself has no model of the device. The pixel crosstalk, the phase quantization and the finite fill factor of a real SLM change the trap
intensities after the hologram is computed, and the experiments absorb the change with camera iterations.

Putting a device model into the solver is not new. Schroff et al.\ model the pixel crosstalk by upscaling each SLM pixel to three by three computational pixels and convolving the phase with a fitted kernel inside a conjugate-gradient minimization of a normalized intensity error, and they generate light potentials, among them a Gaussian spot array, with it~\cite{schroff}. Persson et al.\ modify a hologram generating algorithm to account for pixel crosstalk and report, per their abstract, an improvement of the spot uniformity by more than 100 percent on a device with large crosstalk~\cite{persson}. What the published work leaves open is a measured comparison, run by one code on one task, between a solver that has the device model in its loop and the solver that computes the arrays of the experiments, and the answer to the question that follows from Schroff et al.: whether the gain of such a refinement comes from the device model, from the loss that the conjugate-gradient steps minimize, or from the weights of weighted GS.

What Schroff et al.\ already do is the core of the idea: a crosstalk model inside a conjugate-gradient refinement, with the normalized-error loss, on shaped light potentials and on a Gaussian spot array (their Fig.~5c)~\cite{schroff}. We study the transposition of that approach to point-trap arrays as one method, the device-aware conjugate-gradient refinement. It starts from the weighted GS of the experiments, then takes conjugate-gradient steps on a loss that is evaluated on the device and differentiated through it in closed form, with a target level that follows the current efficiency, and it returns the best admissible iterate. Three elements are the method's own: the efficiency-level loss, the multiplicative weight update of weighted GS (an option, off by default because the ablation refutes it) and the best-iterate safeguard. The normalized-error loss of Schroff et al.\ is a second option of the same optimizer, so that the closest published method is run through the same warm start, device model, optimizer and safeguard and differs only in the loss. The paper reports the outcome as a negative result against that method: none of the three elements improves on it, and what lowers the spread is the device model in the loop, which Schroff et al.\ already use. The contributions are the following.
\begin{enumerate}
\item A measurement that isolates the loss, the weight update, the safeguard and the device model on one task, through one optimizer: the efficiency-level loss and the normalized-error loss give spreads that agree within the seed ranges on a lattice behind a crosstalk device, the normalized-error loss has the lower median on the full device, the safeguard does not act, the weight update raises the spread seven times on the acceptance task, and removing the device model from the loop returns the spread to that of phase-fixed WGS (Sections~\ref{sec:results} and~\ref{sec:ablation}).
\item Proposition~\ref{prop:losses}, which shows why the two losses behave alike for equal traps: they have the same value at the phase where an outer iteration sets its target, and their gradients differ by a push along the gradient of the total light (Section~\ref{sec:theory}).
\item The transposition itself, defined in full: the closed-form gradient through the adjoint of the crosstalk operator, the efficiency-level loss, the safeguard, the algorithm, the parameters and the cost, for traps on FFT bins (Section~\ref{sec:method}).
\item The comparison with phase-fixed weighted GS on the same tasks, with seeds, intervals, the geometries of the studio's benchmark read by the benchmark's win rule and the time per hologram on the device: the gain that the device model brings is large on a lattice behind a crosstalk device, partial on the benchmark's device, and it costs minutes to tens of minutes per hologram against seconds (Sections~\ref{sec:design} and~\ref{sec:results}).
\end{enumerate}
Section~\ref{sec:related} places the method among the published ones, Section~\ref{sec:problem} fixes the model, the device and the metrics, Section~\ref{sec:method} defines the method,
Section~\ref{sec:theory} gives the propositions that explain it, and Sections~\ref{sec:design} to~\ref{sec:ablation} give the experiments. The moving-trap sequences of the same studio, which solve a different problem, are the subject of a companion paper. Section~\ref{sec:repro} lists the commands that rebuild every number.
